\section*{الفصل \ref{ch:g7:geological-time} --- أعمار الأرض}

\begin{solution}{exo:g7:geological-time:1}
(أ) \omterm{def:g7:geological-time:unconformity}{عدم التوافق}؛ (ب) الاحتواء؛ (ج) التراكب؛ (د) التتابع الأحفوري؛ (ه)
التقاطع.
\end{solution}

\begin{solution}{exo:g7:geological-time:2}
(أ) الحجر الرملي الأحمر، في الأسفل، هو الأقدم؛ والحجر الرملي الأصفر، في
الأعلى، هو الأحدث. (ب) \omterm{def:g7:geological-time:superposition}{مبدأ التراكب}. (ج) ألا تكون الطبقات قد انقلبت منذ
تكونت.
\end{solution}

\begin{solution}{exo:g7:geological-time:3}
(أ) حقبة الحياة القديمة، فحقبة الحياة الوسطى، فحقبة الحياة الحديثة. (ب)
حقبة الحياة الحديثة. (ج) حقبة الحياة الوسطى. (د) الجوراسي يأتي مباشرة بعد
الترياسي؛ والطباشيري يأتي مباشرة قبل الباليوجيني.
\end{solution}

\begin{solution}{exo:g7:geological-time:4}
(أ) خطأ: \omterm{def:g7:geological-time:relative-dating}{التأريخ النسبي} لا يفعل إلا ترتيب \omterm{def:g3:stones-pebbles-sand:rock}{الصخور}؛ أما العمر بالسنين فهو \omterm{def:g7:geological-time:relative-dating}{عمر
مطلق}. (ب) صواب. (ج) خطأ: \omterm{def:g3:stones-pebbles-sand:pebble}{الحصى} أقدم من الرصيص الذي يحتويها (\omterm{def:g7:geological-time:inclusion}{مبدأ
الاحتواء}). (د) صواب (\omterm{def:g7:geological-time:fossil-succession}{مبدأ التتابع الأحفوري}).
\end{solution}

\begin{solution}{exo:g7:geological-time:5}
(أ) أحدث: فهو يقطع الطبقة 3، \omterm{def:g3:stones-pebbles-sand:rock}{والصخر} الذي يقطع غيره أحدث منه (التقاطع).
(ب) أقدم: فالطبقة 4 تستقر فوقه دون أن يقطعها. (ج) الطبقة 1، فالطبقة 2،
فالطبقة 3، فالغرانيت، فالطبقة 4. (د) تآكلت \omterm{def:g3:stones-pebbles-sand:rock}{الصخور} التي فوقه بفعل \omterm{def:g5:water-shapes-the-land:erosion}{التعرية}
حتى بلغ الغرانيت السطح؛ ثم وُضعت الطبقة 4 على سطح \omterm{def:g5:water-shapes-the-land:erosion}{التعرية} هذا، وهو \omterm{def:g7:geological-time:unconformity}{عدم
توافق}.
\end{solution}

\begin{solution}{exo:g7:geological-time:6}
(أ) $538.8 - 251.9 = 286.9$ مليون سنة. (ب) $251.9 - 66 = 185.9$
مليون سنة. (ج) 66 مليون سنة. (د) حقبة الحياة القديمة هي الأطول، وحقبة
الحياة الحديثة هي الأقصر (حتى الآن).
\end{solution}

\begin{solution}{exo:g7:geological-time:7}
(أ) $4540 - 538.8 = 4001.2$ مليون سنة. (ب) $4001.2 \div 4540
\approx 0.881$: نحو \textbf{88~\%}. (ج) قبل حقبة الحياة القديمة، لم تكن
الحيوانات ذات الأصداف شائعة بعد، والكائنات الحية الطرية لا تكاد تصير
\omterm{def:g7:sediments-and-fossils:fossil}{أحافير} أبدًا؛ أما الأصداف والعظام، التي تتحفّر جيدًا، فلم تصر شائعة إلا
ابتداء من حقبة الحياة القديمة.
\end{solution}

\begin{solution}{exo:g7:geological-time:8}
(أ) $4540 \div 100 = 45.4$ مليون سنة لكل متر. (ب) $538.8 \div 45.4
\approx 11.9$: نحو \qty{11.9}{m} قبل خط النهاية. (ج) $66 \div
45.4 \approx 1.45$: نحو \qty{1.45}{m}. (د) $0.3 \div 45.4 \approx
0.0066$~متر، أي نحو \qty{0.66}{cm}، أقل من سنتيمتر واحد: لا يكاد العدّاء
يلاحظها.
\end{solution}

\begin{solution}{exo:g7:geological-time:9}
(أ) جاء السيلوري أولًا. نعم: \omterm{def:g3:stones-pebbles-sand:rock}{فالصخر} الأقدم يقع تحت الأحدث. (ب) توضع
الطبقات الرمادية مستوية؛ تُدفع حتى تنتصب؛ تتآكل؛ يوضع \omterm{def:g3:stones-pebbles-sand:sand}{الرمل} الأحمر فوقها.
(ج) لأنه سطح \omterm{def:g5:water-shapes-the-land:erosion}{تعرية} بين \omterm{def:g3:stones-pebbles-sand:rock}{صخور} أعمارها متباعدة جدًا: فالزمن الطويل الذي
ارتفعت فيه الطبقات الرمادية وتآكلت لم يترك وراءه أي طبقة.
\end{solution}

\begin{solution}{exo:g7:geological-time:10}
(أ) السجّيلان اللذان فيهما أمونيت من العمر نفسه، وكذلك الحجران الجيريان
اللذان فيهما قنافذ \omterm{def:g2:water-on-earth:ocean}{بحر}: \omterm{def:g7:sediments-and-fossils:fossil}{الأحافير} نفسها، فالعمر نفسه (\omterm{def:g7:geological-time:fossil-succession}{مبدأ التتابع
الأحفوري}). (ب) الأقدم هو الحجر الرملي الذي فيه ثلاثيات فصوص في الجرف A
(فهو يقع تحت السجّيل الذي فيه أمونيت، وهذا من عمر الطبقة السفلى في الجرف
B). والأحدث هو الطين الذي فيه أوراق نبات في الجرف B. (ج) من الأسفل إلى
الأعلى: حجر رملي فيه ثلاثيات فصوص؛ سجّيل فيه أمونيت؛ حجر جيري فيه قنافذ
\omterm{def:g2:water-on-earth:ocean}{بحر}؛ طين فيه أوراق نبات.
\end{solution}

\begin{solution}{exo:g7:geological-time:11}
(أ) الشست. فهو يقع تحت الرصيص (التراكب)، \omterm{def:g3:stones-pebbles-sand:pebble}{وحصى} الرصيص قطع من الشست
(الاحتواء)؛ وسطح \omterm{def:g5:water-shapes-the-land:erosion}{التعرية} يقطع الشست أيضًا. (ب) بعده: فالحجر الرملي وُضع
مستويًا، إذن مال لاحقًا. (ج) قبله: فالطين يغطي \omterm{def:g7:earthquakes:fault}{الصدع} دون أن ينكسر. (د)
يتكون الشست؛ ويتآكل (السطح المتموج)؛ ويوضع الرصيص والحجر الرملي والحجر
الجيري، بهذا الترتيب؛ وتميل الطبقات، ويكسرها \omterm{def:g7:earthquakes:fault}{الصدع} --- والمقطع وحده لا يقول
أي هذين جاء أولًا؛ وتتآكل الطبقات المائلة \omterm{def:g7:earthquakes:fault}{والصدع} حتى تستوي؛ ويوضع الطين
فوقها. (ه) حالتان: السطح المتموج بين الشست والرصيص، والسطح بين الطبقات
المائلة والطين المستوي.
\end{solution}

\begin{solution}{exo:g7:geological-time:12}
(أ) نصف 365 يومًا هو 182.5 يومًا. والأشهر من يناير إلى يونيو مجموعها
$31+28+31+30+31+30 = 181$ يومًا، فينقضي نصف التاريخ بعد يوم ونصف من بداية
يوليو: في \textbf{2~يوليو عند الظهر}. (ب) $4540 \div 2 = 2270$ مليون
سنة. (ج) $80 \div \num{4540000000} \times \num{31536000} \approx 0.56$:
نحو \textbf{0.6~ثانية}.
\end{solution}

\begin{solution}{pb:g7:geological-time:1}
\textbf{1.} \num{4540} مليون سنة.

\textbf{2.} $4540 \div 365 \approx 12.4$ مليون سنة.

\textbf{3.} $\num{4540000000} \div (365 \times 24) = \num{4540000000}
\div 8760 \approx \num{518000}$ سنة (وبالقيمة المقرَّبة \num{12.4}
مليون، $\num{12400000} \div 24 \approx \num{517000}$).

\textbf{4.} $\num{4540000000} \div \num{525600} \approx 8638$: نحو
\num{8600} سنة.

\textbf{5.} $4540 - 4030 = 510$ مليون سنة بعد تكوّن الأرض؛
$510 \div 12.4 \approx 41$ يومًا. وفي يناير 31 يومًا، فتكون 41 يومًا
قد مضت عند بداية \textbf{11~فبراير}.

\textbf{6.} $4540 - 3500 = 1040$ مليون سنة؛ $1040 \div 12.4 \approx
83.9$ يومًا (83.6 بالقيمة غير المقرَّبة $12.44$). ويناير وفبراير مجموعهما
59 يومًا: أي 24.6 إلى 24.9 يومًا من مارس، في \textbf{25~مارس}.

\textbf{7.} $66 \div 12.4 \approx 5.3$ أيام قبل نهاية السنة،
أي $365 - 5.3 = 359.7$ يومًا بعد 1~يناير. والأشهر الأحد عشر الأولى
مجموعها 334 يومًا: أي 25.7 يومًا من ديسمبر، في \textbf{26~ديسمبر}.

\textbf{8.} $538.8 \div 4540 \approx 0.119$: يبقى \textbf{11.9~\%} من
السنة، أي $0.119 \times 365 \approx 43$ يومًا. و$365 - 43 =
322$ يومًا بعد 1~يناير: فهي تبدأ في \textbf{نوفمبر} (في
18~نوفمبر).

\textbf{9.} $\num{10000} \div 8600 \approx 1.2$ دقيقة.

\textbf{10.} $5000 \div 8600 \approx 0.58$ دقيقة، و$0.58 \times 60
\approx 35$: نحو \textbf{35 ثانية} قبل منتصف الليل.

\textbf{11.} $\num{300000} \div 8600 \approx 34.9$: نحو \textbf{35
دقيقة} قبل منتصف الليل.

\textbf{12.} \babelsublr{35} دقيقة قبل منتصف ليل 31~ديسمبر: يظهر نوعنا نحو الساعة
\textbf{\babelsublr{23:25}}، في الساعة الأخيرة من السنة.
\end{solution}
